\subsection{Experiment Overview}
\label{sec:experiment_overview}

The experiments follow a deliberate logical progression. Experiment~1 first selects the
best Stage-A channel-selection configuration; Experiment~2 then benchmarks that chosen
configuration against competing detectors; and Experiments~3--8 stress-test and
characterise the resulting pipeline along orthogonal axes (epoch duration, boundary
tolerance, fallback sensitivity, blink morphology, epoch-health exclusion, and long-blink
recall). Fixing the channel configuration first, before any baseline comparison, ensures
that the headline comparison in Experiment~2 uses a principled spatial selection rather
than an arbitrary default, and that every subsequent robustness check operates on the same
ideal configuration. Each experiment is summarised below in terms of its purpose,
motivation, and objective.

\paragraph{Experiment~1 --- Channel-selection strategies.}
\textit{Purpose:} determine which Stage-A channel-selection group yields the best detector,
evaluated on both datasets (Raja and Cao2018) and under both estimators (median and mean),
before any other experiment is run.
\textit{Motivation:} the current default rule flags an epoch whenever \emph{any} channel
over \emph{all} channels exceeds its peak-to-peak threshold, which can be triggered by a
single noisy non-ocular channel and inflate the number of flagged epochs; a more selective
spatial group may screen blinks more cleanly.
\textit{Objective:} run the complete Stage~A$\to$B$\to$C pipeline on each candidate
channel group and select the best group by weighing multiple criteria rather than a single
metric. The decision criteria are:
\begin{description}
  \item[(i) Downstream event $F_1$.] The end-to-end event-level $F_1$ of the full pipeline
    when restricted to the group; this is the ultimate quantity the detector is optimised
    for and is therefore the primary criterion.
  \item[(ii) Stage-A epoch-selection $F_1$.] The precision and recall of the
    suspicious-epoch selection itself, scored against epoch-level ground truth (an epoch is
    positive if it contains at least one annotated blink); this isolates the quality of the
    screening stage from the downstream localisation.
  \item[(iii) Number of channels for Stage~A.] The count of channels the group requires,
    as a measure of montage cost, simplicity, and portability; a group that performs
    comparably with fewer channels is preferred because it is cheaper to acquire and easier
    to deploy.
  \item[(iv) False-positive epoch rate / \% epochs flagged.] The fraction of non-blink
    epochs erroneously flagged and the overall proportion of epochs flagged; a low value
    indicates robustness to artefacts and avoids over-flagging that would dilute the
    Stage-B threshold estimate.
  \item[(v) Recall on safety-critical blinks.] The recall specifically on long
    drowsiness-related closures, reflecting a ``do-not-miss'' requirement: a configuration
    that maximises overall $F_1$ but systematically misses safety-critical blinks is
    undesirable.
  \item[(vi) Cross-dataset consistency.] The agreement of the chosen group's behaviour
    between Raja and Cao2018; a group that is best on one dataset but poor on the other does
    not generalise.
  \item[(vii) Best-channel stability.] The stability of which channel (or channels) drives
    detection across sessions and subjects; an unstable best channel signals an
    overfitted, fragile selection.
  \item[(viii) Estimator robustness.] The consistency of the group's ranking under the
    median versus mean centre at Stage~B; a group whose advantage survives both estimators
    is more trustworthy than one that depends on a particular centre.
\end{description}
The configuration that best balances these criteria is designated the ideal configuration
and is carried forward unchanged into Experiments~2--8.

\paragraph{Experiment~2 --- Strategy comparison.}
\textit{Purpose:} benchmark the proposed pipeline, configured with the ideal channel group
from Experiment~1, against competing detection algorithms.
\textit{Motivation:} a method is only convincing if it outperforms the established
alternatives under identical preprocessing, epoching, and scoring.
\textit{Objective:} compare BLINKER-concat, MNE-annot, DBO, Proposed-Mean, and Proposed-Med
at the event level on both datasets, reporting precision, recall, and $F_1$ with paired
significance tests, and isolating both the estimator contrast (median versus mean) and the
cross-dataset generalisation gap.

\paragraph{Experiment~3 --- Epoch-duration stability.}
\textit{Purpose:} characterise how sensitive the proposed pipeline is to the choice of
epoch length.
\textit{Motivation:} the \(30\)~s primary epoch is a design choice, and a robust method
should not depend critically on it.
\textit{Objective:} sweep the epoch duration (e.g.\ \(10/20/40/60\)~s around the \(30\)~s
reference) and quantify the resulting variation in event-level $F_1$, testing each duration
against the reference with Wilcoxon signed-rank tests.

\paragraph{Experiment~4 --- Boundary-tolerance (IoU) stability.}
\textit{Purpose:} assess how the results depend on how strictly a detection must overlap a
ground-truth interval to count as a match.
\textit{Motivation:} blink boundaries are inherently uncertain, so the matching tolerance
is a scoring choice that could materially change the reported numbers.
\textit{Objective:} recompute event-level $F_1$ across a range of intersection-over-union
matching thresholds (e.g.\ \(\{0.0, 0.1, 0.2, 0.3, 0.5\}\)) and report the range of
$F_1$ as a measure of boundary-tolerance sensitivity.

\paragraph{Experiment~5 --- $n_{\min}$ / Stage-B fallback sensitivity.}
\textit{Purpose:} examine the behaviour of the Stage-B fallback that switches from the
flagged epochs to the full valid-epoch set when too few epochs are flagged.
\textit{Motivation:} the fallback count \(n_{\min}\) guards against estimating a threshold
from an insufficient sub-sample, but the choice of \(n_{\min}\) should not unduly drive the
results.
\textit{Objective:} quantify how threshold variance depends on the sub-sample size and how
often the fallback is triggered as a function of epoch duration, characterising the
estimator's stability under small flagged sets.

\paragraph{Experiment~6 --- Blink-region morphology (butterfly).}
\textit{Purpose:} visually and quantitatively characterise the morphology of the detected
blink regions.
\textit{Motivation:} aggregate $F_1$ does not reveal \emph{why} a detection succeeds or
fails; inspecting the waveforms distinguishes genuine, time-locked blink shapes from ragged
artefactual ones.
\textit{Objective:} produce butterfly overlays of true-positive, false-negative, and
false-positive blink regions, stratified by duration and amplitude, per subject and pooled,
to characterise the kinds of events the detector recovers and misses.

\paragraph{Experiment~7 --- Effect of excluding low-health epochs.}
\textit{Purpose:} measure how much the reported performance depends on the preprocessing
decision to exclude low-health epochs, for every detector.
\textit{Motivation:} dropping low-health epochs is a preprocessing step applied to all
methods, not part of any detector, so a reviewer will rightly ask whether the headline
numbers are partly an artefact of that exclusion.
\textit{Objective:} run every condition with and without the exclusion (health-on versus
health-off) on both datasets, reporting the per-condition change in $F_1$ to separate
preprocessing effects from detector effects and to test whether the proposed method depends
on the exclusion less than the baselines.

\paragraph{Experiment~8 --- Long-blink (drowsiness closure) recall.}
\textit{Purpose:} quantify recall specifically on long eye closures for every detector.
\textit{Motivation:} in a drowsy-driving context the long closures (microsleep,
\(\geq 0.5\)~s) are the safety-critical events, yet shape-based detectors tuned to short
blinks tend to reject wide-plateau closures.
\textit{Objective:} split blinks into normal (\(<0.5\)~s) and long (\(\geq 0.5\)~s)
categories and report the per-detector recall gap on both datasets, turning ``the method
also handles long closures'' from a claim into a measured result.
